% 4. Sistem za automatizovanu evaluaciju modela
\section{Sistem za automatizovanu evaluaciju modela}
\label{sec:pipeline}

Za potrebe evaluacije opisane u poglavlju~\ref{sec:metodologija} razvijen je zaseban sistem za automatizovanu evaluaciju modela, organizovan kao DVC \emph{pipeline} (engl. \emph{Data Version Control})\footnote{\url{https://dvc.org}}.
DVC omogućava deklarativno opisivanje faza obrade sa eksplicitnim zavisnostima i izlazima: svaka faza se ponovo izvršava samo ako su se njeni ulazi ili parametri promenili, što eksperimente čini reproducibilnim i štedi troškove API poziva.

\subsection{Format ulaznih podataka}

Svaki od devet skupova podataka (pet kurseva, pri čemu su za pojedine kurseve obuhvaćeni odvojeni ispitni rokovi) opisan je dvema CSV datotekama:

\begin{itemize}
    \item \texttt{questions.csv}, po jedan red za svako pitanje, sa kolonama: identifikator pitanja, tekst pitanja, tačan odgovor nastavnika i adresa nastavnog materijala (PDF udžbenika) iz koga je pitanje postavljeno;
    \item \texttt{student\_answers.csv}, po jedan red za svaki studentski odgovor, sa kolonama: identifikator studenta, identifikator pitanja, ocena nastavnika (0--10) i tekst odgovora.
\end{itemize}

\subsection{Faze obrade}

Prva faza (\emph{spajanje ulaza}) validira i spaja ove dve datoteke u jedinstvenu tabelu.
Validacija obuhvata proveru obaveznih kolona, jedinstvenosti identifikatora pitanja, nepraznosti svih polja i kompletnosti podataka, svaki student mora imati odgovor na svako pitanje; u suprotnom se obrada prekida sa opisom greške.
Ovakve stroge provere pokazale su se ključnim za kvalitet eksperimenata, jer su greške u ručno pripremanim podacima otkrivane pre nego što bi izazvale skupe, delimično neispravne eksperimente.

Druga faza (\emph{ocenjivanje}) za svaki red spojene tabele konstruiše prompt i poziva izabrani jezički model preko integracionog sloja iz odeljka~\ref{sec:integracije}.
Postoje tri varijante ove faze, od kojih dve odgovaraju režimima ocenjivanja iz odeljka~\ref{sec:modeli}:

\begin{itemize}
    \item \textbf{PRA varijanta} u prompt uključuje pitanje, tačan odgovor nastavnika i studentski odgovor.
    \item \textbf{GRA varijanta} umesto tačnog odgovora nastavnika koristi referentni odgovor koji model sam generiše iz nastavnog materijala.
    Ovo je dvostepen postupak: u prvom koraku materijal se preuzima sa zadate adrese kao PDF, tekst se izdvaja bibliotekom PyPDF2 i kešira na disku (po heš vrednosti adrese, tako da se svaki udžbenik preuzima i obrađuje samo jednom), a model iz njega konstruiše koncizan referentni odgovor dužine dve do pet rečenica.
    U drugom koraku studentski odgovor ocenjuje se u odnosu na taj generisani standard, istim postupkom kao u PRA varijanti.
    Time se za svaki skup podataka dobija tabela pitanja u kojoj su nastavnički referentni odgovori zamenjeni generisanim, pa se ocenjivanje sprovodi ponovnom upotrebom iste komponente kao u PRA varijanti.
    \item \textbf{Varijanta sa nastavnim materijalom u promptu} ne koristi referentni odgovor, već ceo izvučeni tekst udžbenika umeće u svaki prompt za ocenjivanje.
    Ova varijanta razmatrana je, ali nije sprovedena do kraja i njeni rezultati se ne koriste u analizi; razlozi su izloženi u odeljku~\ref{sec:rezim-udzbenik}.
\end{itemize}

U svim varijantama sistemska instrukcija se sastoji od prompta strogosti (blag, neutralan ili strog; dodatak~\ref{app:promptovi}) i fiksnog uputstva o formatu izlaza: prvi red odgovora mora sadržati samo ocenu na skali 0--100, a ostatak obrazloženje.
Temperatura se ne postavlja eksplicitno, pa modeli rade na podrazumevanoj vrednosti svog provajdera.
Numerička ocena se parsira iz prvog reda odgovora; skala 0--100 naknadno se preslikava na skalu 0--10 radi poređenja sa ocenama nastavnika.
Ako parsiranje ne uspe, poziv se ponavlja, do tri pokušaja ukupno; ako API poziv privremeno otkaže, ponavljanju prethodi rastuće odlaganje.
Trajno neuspešni redovi označavaju se specijalnom negativnom vrednošću i isključuju iz analize.

\subsection{Automatizacija eksperimenata i analiza}

Potpuni eksperiment obuhvata petnaest kombinacija (pet modela $\times$ tri nivoa strogosti) nad svih devet skupova podataka, dakle 135 pokretanja po režimu ocenjivanja.
Pomoćna skripta automatizuje ovaj proces: za svaku kombinaciju upisuje izbor modela i strogosti u parametarsku datoteku i ponovo pokreće DVC \emph{pipeline}, pri čemu ime svake izlazne datoteke kodira konfiguraciju (strogost, model i režim) kojom je proizvedena.
Svaka kombinacija pokrenuta je jednom, pa rezultati ne obuhvataju procenu varijabilnosti između ponovljenih pokretanja istog modela.

Završna analitička skripta prolazi kroz sve ocenjene datoteke, uparuje ocene modela sa ocenama nastavnika i računa metrike opisane u odeljku~\ref{sec:metrike}, Pirsonov koeficijent korelacije, odnos prosečnih ocena, RMSE i odnos RMSE/STDDEV, po svakoj kombinaciji skupa podataka, modela, strogosti i režima.
Rezultati se objedinjuju u zbirnu tabelu iz koje su izvedeni rezultati prikazani u poglavlju~\ref{sec:rezultati}.
